\section{An involution with a prescribed perturbation}
\label{inv:section}
\label{sec:involution}

The construction uses a Hamiltonian involution whose fixed components
are smaller than the ambient manifold. A small invariant perturbation
turns a chosen Morse function on each component into an exact list of
nondegenerate Hamiltonian fixed points. The following elementary estimate
excludes additional fixed points. We give the involution argument of
\cite[Proposition~4.2]{OpenAI2026Stable} in a general form; the period estimate is a
compact-manifold version of the classical Lipschitz period bound
\cite{Yorke1969}.

\begin{lemma}\label{inv:short-period}
Let $V$ be a smooth vector field on a closed manifold $P$, with flow
$\psi^s$. There exists $\tau>0$ such that
\[
  0<T<\tau,\qquad \psi^T(x)=x
  \quad\Longrightarrow\quad V(x)=0.
\]
\end{lemma}

\begin{proof}
Choose finitely many coordinate balls, each relatively compact in a
larger coordinate chart, and smaller balls that cover $P$. Uniform speed
bounds give $\delta>0$ such that a trajectory starting in a smaller ball
stays in its corresponding coordinate ball for time $\delta$. Choose
a common Lipschitz constant $L>0$ for the coordinate vector fields on
these balls.

Suppose $\psi^T(x)=x$ with $T<\delta$. Write the trajectory in one of
these charts as $y:[0,T]\to\R^d$, with $\dot y=v(y)$. Set
\[
  m=\max_{0\leq t\leq T}|v(y(t))|
\]
and choose $t_0$ realizing the maximum. Since $y(T)=y(0)$,
$\int_0^T v(y(t))\,dt=0$. Moreover,
$|y(t)-y(t_0)|\leq Tm$. Consequently
\[
\begin{split}
  m
  &=\left|\frac1T\int_0^T
       \bigl(v(y(t_0))-v(y(t))\bigr)\,dt\right|\\
  &\leq\frac LT\int_0^T|y(t_0)-y(t)|\,dt
   \leq LTm.
\end{split}
\]
For $T<\min\{\delta,L^{-1}\}$ this forces $m=0$.
\end{proof}

We use the convention $\omega(X_f,\cdot)=df$ for Hamiltonian vector
fields. A fixed point of a diffeomorphism $\Phi$ is nondegenerate if
$d\Phi-I$ is invertible there.

\begin{proposition}\label{inv:perturb}
Let $(P,\omega)$ be a closed symplectic manifold and let $g$ be a
Hamiltonian diffeomorphism with $g^2=\mathrm{id}$. Its fixed locus is a
finite disjoint union of closed symplectic submanifolds,
\[
  \Fix(g)=F_1\sqcup\cdots\sqcup F_r.
\]
For each $j$, prescribe a smooth Morse function $h_j:F_j\to\R$.
There is a smooth $g$-invariant function $f:P\to\R$ restricting to $h_j$
on $F_j$ such that, for every sufficiently small $\epsilon>0$,
\[
  \Fix(\psi^\epsilon\circ g)
  =\bigsqcup_{j=1}^r\operatorname{Crit}(h_j),
\]
where $\psi^s$ is the Hamiltonian flow of $f$. Every fixed point in this
list is nondegenerate. The map $\psi^\epsilon\circ g$ is the time-one map
of a smooth one-periodic Hamiltonian. If $P$ is simply connected, all
fixed-point loops for that Hamiltonian are contractible.
\end{proposition}

\begin{proof}
Average a Riemannian metric over the involution. At a fixed point $x$,
the exponential map for this metric intertwines $g$ with $d_xg$.
Thus the fixed locus is a smooth closed submanifold whose tangent space
at $x$ is
\[
  E_x^+=\ker(d_xg-I).
\]
The local description and compactness give finitely many connected
components. Since $(d_xg)^2=I$, there is a direct-sum decomposition
\[
  T_xP=E_x^+\oplus E_x^- ,\qquad
  E_x^-=\ker(d_xg+I).
\]
For $u\in E_x^+$ and $v\in E_x^-$, symplectic invariance gives
$\omega_x(u,v)=\omega_x(d_xg(u),d_xg(v))=-\omega_x(u,v)$.
The summands are therefore symplectically orthogonal, and the restriction
of $\omega_x$ to each is nondegenerate. In particular, each $F_j$ is
symplectic and $E_x^-$ is its symplectic normal space.

Extend the $h_j$ simultaneously to a smooth function $\widetilde f$ on
$P$, using disjoint tubular neighborhoods of the components and cutoff
functions. Set
\[
  f=\tfrac12(\widetilde f+\widetilde f\circ g).
\]
Then $f\circ g=f$ and $f|_{F_j}=h_j$. Because $g$ is symplectic,
invariance implies $g_*X_f=X_f$, so $g$ commutes with $\psi^s$. At
$x\in F_j$, invariance also gives $d_xf(v)=0$ for $v\in E_x^-$.
It follows that $X_f(x)\in T_xF_j$ and that $X_f|_{F_j}$ is the
Hamiltonian vector field of $h_j$ for $\omega|_{F_j}$. In particular,
\begin{equation}\label{inv:equilibria}
  x\in F_j,\quad X_f(x)=0
  \quad\Longleftrightarrow\quad d_xh_j=0.
\end{equation}

Write $\Phi_\epsilon=\psi^\epsilon\circ g$. Commutation gives
\[
  \Phi_\epsilon^2=\psi^{2\epsilon}.
\]
Choose $2\epsilon<\tau$, where $\tau$ is supplied by
Lemma~\ref{inv:short-period} for $X_f$. If $\Phi_\epsilon(x)=x$, then
$\psi^{2\epsilon}(x)=x$, so $X_f(x)=0$. Thus $\psi^\epsilon(x)=x$,
and the original fixed-point equation implies $g(x)=x$.
Equation~\eqref{inv:equilibria} identifies $x$ as a critical point of
one of the $h_j$. Conversely, every such critical point is fixed by
both $g$ and $\psi^\epsilon$. This proves the asserted equality of fixed
sets.

We next check nondegeneracy in the entire tangent space. At one of these
equilibria let $A_x$ be the linearization of $X_f$. Then
$d_x\psi^\epsilon=\exp(\epsilon A_x)$. Equivariance implies that
$A_x$ preserves $E_x^+$ and $E_x^-$; denote its restrictions by
$A_x^+$ and $A_x^-$. On $E_x^+=T_xF_j$, the defining identity for
$X_f$ gives
\[
  \omega_x(A_x^+u,v)=\operatorname{Hess}_x(h_j)(u,v).
\]
Thus $A_x^+$ is invertible. The limits
\[
  \frac{\exp(\epsilon A_x^+)-I}{\epsilon}\longrightarrow A_x^+,
  \qquad
  -\exp(\epsilon A_x^-)-I\longrightarrow-2I
\]
show that both blocks of $d_x\Phi_\epsilon-I$ are invertible for all
sufficiently small positive $\epsilon$. A single choice of $\epsilon$
works at the finite set of critical points.

Finally, let $K_s$, $0\leq s\leq1$, be a smooth Hamiltonian whose
time-one map is $g$. Choose nonnegative smooth functions $a,b$, each
of integral one, supported in successive disjoint subintervals of
$(0,1)$, and put
\[
  A(t)=\int_0^t a(u)\,du,\qquad
  H_\epsilon(t,x)=a(t)K_{A(t)}(x)+\epsilon b(t)f(x).
\]
The first interval traverses the Hamiltonian path of $K_s$, and the
second traverses $\psi^s$ from $s=0$ to $s=\epsilon$. Hence the time-one
map is $\psi^\epsilon\circ g$. Since $H_\epsilon$ vanishes near $t=0$
and $t=1$, its periodic extension is smooth. On a simply connected $P$,
every loop traced by a fixed point along this path is contractible.
\end{proof}

The surface below supplies four fixed components and an exceptional
sphere on which the first Chern class evaluates to one. It is the
surface of \cite[Proposition~4.3]{OpenAI2026Stable}.

\begin{proposition}\label{inv:surface}
Let $Y$ be the complex blow-up of $\CP^1\times\CP^1$ at
$\{0,\infty\}\times\{0,\infty\}$. There is a K\"ahler form
$\omega_Y$ for which the lifted diagonal circle action is Hamiltonian.
Its half-turn $g_Y$ fixes exactly the four exceptional curves
$E_1,\ldots,E_4$, each isomorphic to $\CP^1$. Along each curve its
derivative is the identity tangentially and minus the identity on the
symplectic normal bundle. The surface is simply connected, with
\[
  H_i(Y;\Z)\cong
  \begin{cases}
    \Z,&i=0,4,\\
    \Z^6,&i=2,\\
    0,&\text{otherwise},
  \end{cases}
  \qquad
  \langle c_1(TY),[E_j]\rangle=1.
\]
\end{proposition}

\begin{proof}
The blow-up is a smooth projective surface and therefore admits a
K\"ahler form. On affine coordinates the diagonal action is
\[
  t\cdot(z,w)=(e^{2\pi it}z,e^{2\pi it}w),
  \qquad t\in\R/\Z.
\]
It preserves the four centers and lifts holomorphically to $Y$.
Averaging a K\"ahler form over this action gives an invariant
K\"ahler form $\omega_Y$.

Removing four points from the simply connected real four-manifold
$\CP^1\times\CP^1$ preserves simple connectivity: paths and their
contracting disks can be moved off the points by general position.
Every loop in $Y$ can similarly be moved off the exceptional curves,
which have real codimension two. The complement is identified by
blow-down with the punctured original surface, so the loop contracts.
Hence $Y$ is simply connected. If $V$ is the infinitesimal generator
of the lifted circle action, Cartan's formula yields
\[
  d\bigl(\omega_Y(V,\cdot)\bigr)
  =\mathcal L_V\omega_Y=0.
\]
Since $H^1(Y;\R)=0$, there is a smooth $k:Y\to\R$ with
$dk=\omega_Y(V,\cdot)$. The circle action is therefore Hamiltonian,
and the autonomous Hamiltonian $k/2$ generates its half-turn at time one.

Downstairs, the half-turn fixes exactly the four centers. At any center
the circle weights are $(\alpha,\beta)$ with
$\alpha,\beta\in\{1,-1\}$, so the half-turn is
$(z_1,z_2)\mapsto(-z_1,-z_2)$. The local blow-up consists of pairs
$(\ell,v)$ with $\ell\subset\C^2$ a complex line and $v\in\ell$.
The lifted half-turn is
\[
  (\ell,v)\longmapsto(\ell,-v).
\]
It fixes exactly the exceptional zero section, acts identically on
its tangent bundle, and acts by $-1$ on its normal line. Equivariant
blow-down excludes any further fixed points. The symplectic normal
description follows from the eigenspace calculation in
Proposition~\ref{inv:perturb}.

Each point blow-up replaces a four-ball by a disk bundle over $\CP^1$
with common boundary $S^3$. Mayer--Vietoris therefore adds one free
generator to $H_2$ and leaves $H_1$ unchanged. Together with simple
connectivity and integral Poincar\'e duality, this gives the displayed
homology groups from the ranks $(1,2,1)$ of
$\CP^1\times\CP^1$. Finally, the normal bundle of an exceptional
curve is $\mathcal O_{\CP^1}(-1)$. The tangent-normal exact sequence
gives
\[
  \langle c_1(TY),[E_j]\rangle
  =\deg(T\CP^1)+\deg\mathcal O_{\CP^1}(-1)=2-1=1.
\]
\end{proof}

\begin{corollary}\label{inv:realization}
Let $(X,\omega_X)$ be a simply connected closed K\"ahler manifold.
Set $M=X\times Y$, with the product form, and $C=X\times\CP^1$.
For every smooth Morse function $h:C\to\R$, there is a smooth
one-periodic Hamiltonian $H$ on $M$ whose fixed points are all
nondegenerate and satisfy
\[
  \#\Fix(\phi_H^1)=\#\Fix_0(\phi_H^1;H)
  =4\,\#\operatorname{Crit}(h).
\]
Moreover, $M$ is simply connected and K\"ahler, and its minimal Chern
number is one.
\end{corollary}

\begin{proof}
The Hamiltonian involution $g=\mathrm{id}_X\times g_Y$ has precisely
the four fixed components $X\times E_j$, each diffeomorphic to $C$.
Transport $h$ to these components and apply
Proposition~\ref{inv:perturb}. The product $M$ is simply connected,
so all fixed-point loops are contractible. On the sphere
$\{x\}\times E_j$, the $X$-directions form a trivial bundle and
Proposition~\ref{inv:surface} gives
\[
  \langle c_1(TM),[\{x\}\times E_j]\rangle=1.
\]
Thus $c_1(TM)(\pi_2(M))=\Z$.
\end{proof}
